\subsection{Pflichtenheft}
\label{app:Pflichtenheft}

Das Pflichtenheft konkretisiert die technische Umsetzung der im Lastenheft
beschriebenen Anforderungen f"ur das \Fachbegriff{AntispamAdminCenter}.

\subsubsection*{System"ubersicht}
Die Anwendung wird als serverseitige Webanwendung mit Python~3.12, FastAPI,
SQLModel und SQLite umgesetzt. Die Fachlogik ist in Service-Module ausgelagert,
die Oberfl"ache wird "uber HTML-Templates bereitgestellt. Der Zugriff erfolgt
per Browser im internen Netzwerk.

\subsubsection*{Umzusetzende Funktionen}
\begin{enumerate}
\item \textbf{Firmenverwaltung}\\
      Firmen werden in einer zentralen Datenbank gespeichert und "uber
      eigene Formulare erstellt, bearbeitet und angezeigt.
\item \textbf{Benutzerverwaltung}\\
      Benutzer werden mit Vorname, Nachname, E-Mail-Adresse, Rolle und optionaler
      Firmenzuordnung verwaltet. Root- und Admin-Benutzer k"onnen
      Benutzer anlegen und "andern.
\item \textbf{Rechtekonzept}\\
      Das Rollenmodell unterscheidet mindestens zwischen Root, Admin und User.
      Root darf firmen"ubergreifend arbeiten. Admin und User sind auf
      zugeordnete Firmen und die dort zul"assigen Funktionen begrenzt.
\item \textbf{Mail-Adressverwaltung}\\
      Mail-Adressen werden firmenbezogen gespeichert, in Formularen
      validiert und "uber die Oberfl"ache pflegbar gemacht.
\item \textbf{IP-Adressverwaltung}\\
      IP-Adressen werden firmenbezogen gespeichert, auf syntaktische
      Korrektheit gepr"uft und "uber die Oberfl"ache pflegbar gemacht.
\item \textbf{Authentifizierung}\\
      Die Anmeldung erfolgt "uber E-Mail-Adresse und Passwort. Passw"orter werden
      mit Argon2 gehasht. Nach erfolgreicher Anmeldung wird eine Session erzeugt.
\item \textbf{Login-Schutz}\\
      Fehlgeschlagene Anmeldungen werden protokolliert. Zus"atzlich werden
      Rate-Limiting und kontobezogene Sperrmechanismen gegen wiederholte
      Fehlanmeldungen umgesetzt.
\item \textbf{Audit-Log}\\
      Erstellen, "Andern und L"oschen fachlicher Datens"atze werden mit
      Benutzerbezug, Aktion, Objektart und "Anderungsdaten protokolliert.
      Sensible Felder wie Passw"orter werden im Audit-Log maskiert.
\item \textbf{Rollback-Unterst"utzung}\\
      F"ur ausgew"ahlte Audit-Eintr"age werden die erforderlichen Informationen
      f"ur ein kontrolliertes Wiederherstellen des vorherigen Zustands
      gespeichert.
\item \textbf{Datenbereitstellung}\\
      Die Anwendung stellt die verwalteten Fachobjekte sowohl "uber
      browserbasierte Oberfl"achen als auch "uber schlanke API-Endpunkte f"ur
      nachgelagerte Verarbeitung bereit.
\end{enumerate}

\subsubsection*{Qualit"atsanforderungen}
\begin{enumerate}
\item Die Fachlogik muss vom Pr"asentationslayer getrennt implementiert werden.
\item Eingaben m"ussen serverseitig validiert werden.
\item Sicherheitsrelevante Vorg"ange m"ussen nachvollziehbar protokolliert
      werden.
\item Die wichtigsten Funktionen m"ussen durch automatisierte Tests abgesichert werden.
\item Die Konfiguration soll "uber zentrale Konfigurationswerte erfolgen.
\end{enumerate}

\subsubsection*{Betriebsbedingungen}
\begin{enumerate}
\item Einsatz im internen Netzwerk ohne "offentlichen Internetzugriff
\item Bereitstellung auf einem System mit Python-Laufzeitumgebung
\item Persistente Datenspeicherung in einer SQLite-Datenbank
\item Nutzung durch interne Administratoren und autorisierte
      Firmenbenutzer
\end{enumerate}
